% =====================================================================
\section{Related Work}
\label{sec:related}
% =====================================================================

\paragraph{Content-addressing in software and programming languages.}
Content-addressed storage is well-established in software systems
(Git~\cite{git}, IPFS~\cite{benet2014ipfs},
Nix~\cite{dolstra2004nix}).  The closest analog in programming
languages is Unison~\cite{unison}, which was designed from the ground
up around content-addressing: names are metadata, refactoring never
invalidates hashes, and the full syntax tree is hashed.  We differ in
three ways: we \emph{retrofit} content-addressing onto an existing
language not designed for it; we hash only the \emph{type}, not the
implementation (justified by proof irrelevance); and canonicalization
in a dependently typed setting is complicated by definitional equality
and universe polymorphism.
\simone{Deepen Unison comparison. Unison: designed around CA from start, hashes full AST, 512-bit SHA3. Ours: retrofit onto Lean, hash type only (proof irrelevance). Make structural difference explicit.}

\paragraph{Formal mathematics infrastructure.}
Hurd's OpenTheory~\cite{hurd2011opentheory} tackled the same core
problem---prover-independent theorem identity---for the HOL family,
using theory interpretations rather than cryptographic hashing.
\simone{Add OpenTheory (Joe Hurd, 2009--2014)---cross-prover package management for HOL family. Closest prior system. Doesn't use crypto hashing but same core problem. CRITICAL.}
Gauthier et al.~\cite{gauthier2019aligning} and M\"uller et
al.~\cite{muller2017alignment} align concepts \emph{across} type
theories; our content-addressing operates \emph{within} one, where
canonical forms can be compared syntactically---a natural prerequisite
for cross-system alignment.
\simone{Deepen alignment work discussion. Gauthier-Kaliszyk solve harder problem across foundations (HOL4/Coq/Mizar). Ours cleaner but narrower (single type theory). Explain intra-system CA as prerequisite for cross-system.}  The OMDoc/MMT
framework~\cite{rabe2013scalable, kohlhase2006omdoc} and the vision of
Kohlhase and Rabe~\cite{kohlhase2017qed} articulate a pluralistic
formal library; we provide a concrete identity mechanism for it.
Existing library management (Mathlib's monorepo~\cite{mathlib4}, Coq's
opam packages~\cite{coq}, Isabelle's AFP~\cite{isabelle_afp}) relies
on name-based identity throughout.

\sam{Discuss HashMath: independent CIC + content addressing + P2P. Complementary---HashMath from scratch, ours within Lean.}

\paragraph{Proof repair and deduplication.}
Ringer et al.~\cite{ringer2021repair} and Gopinathan et
al.~\cite{gopinathan2023repair} repair proofs broken by library
refactoring; content-addressing aims to prevent such breakage rather
than mitigate it.
\simone{Add proof repair literature: Ringer et al.\ PLDI 2021, Gopinathan et al.\ PLDI 2023. Complementary approach---automated repair vs stable identifiers.}  Type-based deduplication of theorem statements is
used in AI-for-theorem-proving pipelines~\cite{syntheticthm},
essentially Level~0 without cryptographic persistence.

\leo{Discuss Yatima project (github.com/argumentcomputer/yatima) as prior work. Dead but had nice ideas---content addressing + zkSNARK for type-checking.}

The question of
normal forms in dependent type theory~\cite{abel2018decidability,
coquand1991algorithm} underlies our canonicalization; we sidestep full
normalization with pragmatic partial levels.

% =====================================================================
